\section{The original normal-crossing carrier with accumulated boundary}
\label{sec:hironaka-normal-crossings-boundary-transform}

\begin{definition}[A simultaneous original carrier and boundary]
\label{hir:general-nc-simultaneous-original-boundary}
\lean{Hironaka.HironakaNormalCrossingsBoundaryTransform.IsBoundaryAdaptedCarrier}
\uses{hir:general-normal-crossings-carrier}
\leanok
A closed original normal-crossing carrier $H$ and a finite list $E$ of
closed hypersurfaces are boundary adapted if, at every point $x$, there
are actual centred holomorphic coordinates on an open neighbourhood,
a finite selected coordinate set $S$, and coordinate sets $B_i$ for the
members of $E$. On the whole neighbourhood, $H$ is the union of the
coordinate hyperplanes indexed by $S$. For each $E_i$ through $x$, $B_i$
is a singleton and $E_i$ is its coordinate hyperplane on that same
neighbourhood. The singletons for the through members are distinct and
disjoint from $S$. There is no condition on the coordinate sets of
members not through $x$. In particular this definition is meaningful
in dimension zero and for empty through families.
\end{definition}

\begin{lemma}[The actual empty-boundary producer and geometric consumers]
\label{hir:general-nc-boundary-empty-producer}
\lean{Hironaka.HironakaNormalCrossingsBoundaryTransform.boundaryAdapted_nil,Hironaka.HironakaNormalCrossingsBoundaryTransform.isNormalCrossingsCarrier_of_boundaryAdapted,Hironaka.HironakaNormalCrossingsBoundaryTransform.isSNCDivisor_of_boundaryAdapted,Hironaka.HironakaNormalCrossingsBoundaryTransform.hasNormalCrossings_highest_boundary}
\uses{hir:general-nc-simultaneous-original-boundary,hir:general-nc-intrinsic-branch-multiplicity}
\leanok
Every actual original NC carrier is boundary adapted with the empty
list. Any boundary adapted pair produces an actual original NC carrier
and an actual SNC boundary. Its intrinsic highest original-branch
stratum has only normal crossings with that actual boundary.
\end{lemma}
\begin{proof}
For the empty list use the actual original coordinates and selected
set; there are no boundary equations or distinctness conditions. For a
general pair, discarding the boundary equations gives an actual NC
flattening for the entire carrier. Discarding the union equation gives
one common actual coordinate system for every through boundary member;
their singleton equations and distinctness give precisely SNC.
At a point of the highest original stratum, the intrinsic count theorem
identifies the whole stratum on the original chart with the intersection
of all selected hyperplanes. Those very equations, together with the
existing singleton boundary equations on the same domain, give actual
normal-crossing coordinates for the whole centre and the boundary.
No centre coordinate system is supplied separately.
\end{proof}

\begin{lemma}[Joint coordinates adapted to the whole highest stratum]
\label{hir:general-nc-highest-original-boundary-coordinates}
\lean{Hironaka.HironakaNormalCrossingsBoundaryTransform.exists_highest_joint_coordinates}
\uses{hir:general-nc-boundary-empty-producer,hir:general-nc-highest-multiplicity-smooth-centre}
\leanok
Suppose all original branch counts are at most $r$. At a point of its
highest stratum there are actual centred coordinates in which the
whole centre is the dimension-$(n-r)$ plane, the entire original
carrier is the union of the last $r$ coordinate hyperplanes, and every
through old boundary member is one of the distinct first $n-r$
coordinate hyperplanes. These coordinates and boundary indices are
internally produced from the simultaneous original pair.
\end{lemma}
\begin{proof}
Take actual simultaneous coordinates and their selected set $S$.
The intrinsic highest-stratum dictionary gives $|S|=r$ and the centre
as the whole intersection of those selected equations. Choose a
permutation sending $S$ onto the last $r$ indices and its complement
onto the first $n-r$ indices. Postcomposition with that holomorphic
linear isomorphism, and precomposition of the actual inverse with its
inverse, gives an actual coordinate map on the same open source.
Transport each old boundary coordinate set by this permutation.
Singletons and their distinctness are preserved. Disjointness from $S$
places every through old boundary singleton in the first block.
The original union and the centre intersection are transported on
the whole actual domain, producing all claimed dictionaries together.
\end{proof}

\begin{lemma}[Actual joint transformed coordinates over the centre]
\label{hir:general-nc-joint-boundary-coordinates-over-centre}
\lean{Hironaka.HironakaNormalCrossingsBoundaryTransform.exists_joint_coordinates_over_highest}
\uses{hir:general-nc-highest-original-boundary-coordinates,hir:general-nc-highest-stratum-local-strict-branch-drop}
\leanok
In the genuine blowup of the highest original stratum, every point
above the centre has one actual centred coordinate system simultaneously
expressing the whole original strict carrier, every through strict
old boundary member and the new exceptional member. The original
selected set has at most $r-1$ indices. All boundary equations are
distinct singleton coordinates disjoint from that selected set.
\end{lemma}
\begin{proof}
Use the produced joint base coordinates. The genuine blowup supplies
a monomial local inverse $W$, pivot $j$ in the last block, and its
actual holomorphic lift inverse. Differentiating those inverse
identities gives an invertible differential for $W$, so its actual
recentred local coordinates exist. In the uncentred monomial chart,
the original strict carrier is exactly the union of the last-block
hyperplanes other than the pivot, by the proved genuine finite-union
strict-closure calculation. A through old boundary member has index
in the first block. Its pivot is distinct, and its original coordinate
is unchanged by the monomial projection. The actual strict-transform
calculation therefore gives exactly that same coordinate hyperplane.
The new exceptional member is exactly the pivot hyperplane.

Every through strict old member maps to a through old base member,
because it lies in the preimage of that closed member. Its coordinate
thus vanishes at the blowup point. The pivot also vanishes because
the point is above the centre. Subtracting the centre value preserves
all these equations. Retain among the remaining original indices
exactly those whose coordinate vanishes at the point, and shrink to
one open neighbourhood on which the finitely many other coordinates
remain nonzero. The entire original union is then precisely the union
of this retained set, whose size is at most $r-1$. All old boundary
indices are tangential and mutually distinct; the new pivot is normal
and is omitted from the retained original set. Hence every through
boundary singleton is distinct and disjoint from the entire retained
original set on the same actual coordinate domain.
\end{proof}

\begin{theorem}[A genuine boundary-preserving original-count step]
\label{hir:general-nc-boundary-transform-original-count-step}
\lean{Hironaka.HironakaNormalCrossingsBoundaryTransform.boundaryAdapted_strict_highest,Hironaka.HironakaNormalCrossingsBoundaryTransform.isSNCDivisor_strict_highest_boundary,Hironaka.HironakaNormalCrossingsBoundaryTransform.branchMultiplicity_boundary_strict_highest_le,Hironaka.HironakaNormalCrossingsBoundaryTransform.exists_highest_boundary_blowup_count_decrease}
\uses{hir:general-nc-joint-boundary-coordinates-over-centre,hir:general-nc-highest-stratum-global-strict-count-decrease,hir:general-nc-boundary-empty-producer}
\leanok
The genuine highest-original-stratum blowup of a boundary adapted pair
produces a boundary adapted pair consisting of the whole strict
original carrier and the list of strict old boundary members followed
by the actual new exceptional member. The new list is genuinely SNC
and every original branch count is at most $r-1$. For a nonempty
original carrier, the positive attained maximum and the entire step
are internally produced without compactness or global original labels.
\end{theorem}
\begin{proof}
All transformed original and old boundary sets are closed by their
strict-closure definition, and the exceptional member is closed by
continuity of the actual projection and closedness of the actual
centre. Above the centre use the jointly produced coordinates of the
preceding lemma. Off the centre use the actual blowup local
biholomorphism to transport one original simultaneous coordinate
system. On its actual open inverse domain, strict transforms of the
closed original and old boundary sets equal their preimages. Thus the
same original selected set and the same old coordinate singletons
give the entire strict original carrier and every through strict old
member together. The exceptional member does not pass through the
point, so assign its unused coordinate set the empty set. Through
singleton cardinality, distinctness and disjointness are inherited
from the actual original pair. These two cases give a boundary adapted
pair at every actual point; its boundary consumer proves SNC.

Above the centre the genuine original-count chart has at most $r-1$
indices. Off it the intrinsic count is unchanged and the image point
is outside the highest stratum, so its integer count is strictly
below $r$. This proves the global bound. If the original carrier is
nonempty, the actual finite intrinsic count image first produces its
positive attained maximum and smooth whole centre, and the proved
complex blowup existence theorem produces the genuine blowup. Applying
the preceding result to these actual objects produces the whole step.
\end{proof}

\begin{remark}
These are statements about actual simultaneous geometric carriers and
boundary equations. They do not assert that the total branch count
decreases, that a separately marked ideal contains the centre in its
cosupport, or that the old ordered monomial algorithm is independent
of its colexicographic branch-label tie-break. Canonical marked-ideal
global sequence descent remains a separate obligation.
\end{remark}
